% !TeX encoding = UTF-8
% !TeX program = pdflatex
% Formatação baseada em atividade2/deliverables/relatorio.tex.
% Arquivo independente: gráficos e diagrama vetoriais incorporados, sem .bib.
% Compilar duas vezes com pdfLaTeX, ou usar Tectonic.
% Materiais da Atividade 3: equipe, evidências, vídeo e referência do repositório.
\documentclass[12pt,a4paper]{article}
\usepackage[T1]{fontenc}
\usepackage[utf8]{inputenc}
\usepackage[brazil]{babel}
\usepackage{mathptmx}
\usepackage{amsmath}
\usepackage{geometry}
\geometry{top=3cm,left=3cm,right=2cm,bottom=2cm}
\usepackage{setspace}
\onehalfspacing
\usepackage{microtype}
\usepackage{booktabs,tabularx,array}
\usepackage{graphicx,float,caption,xcolor}
\usepackage{enumitem,fancyhdr,titlesec}
\usepackage{pgfplots}
\pgfplotsset{compat=1.18}
\usepackage{hyperref}
\usepackage{url}
\usepackage{xurl}
\hypersetup{colorlinks=true,linkcolor=black,citecolor=black,urlcolor=blue,
 pdftitle={Virtualização, namespaces e cgroups: laboratório de contêineres Linux},
 pdfauthor={Equipe GenResearch - UFS}}
\urlstyle{same}
\setlength{\parindent}{1.5cm}
\setlength{\parskip}{0pt}
\setlength{\emergencystretch}{2em}
\setlist{nosep,leftmargin=1.5cm}
\captionsetup{font=small,labelfont=bf,justification=justified,singlelinecheck=false}
\titleformat{\section}[block]{\normalfont\bfseries}{\thesection}{1em}{\MakeUppercase}
\titleformat{\subsection}[block]{\normalfont\mdseries}{\thesubsection}{1em}{\MakeUppercase}
\titlespacing*{\section}{0pt}{0.6\baselineskip}{0.6\baselineskip}
\titlespacing*{\subsection}{0pt}{0.5\baselineskip}{0.3\baselineskip}
\pagestyle{fancy}
\setlength{\headheight}{14pt}
\fancyhf{}
\fancyhead[R]{\small\thepage}
\renewcommand{\headrulewidth}{0pt}
\newcommand{\secao}[1]{\clearpage\section{#1}}
\newcommand{\fontepropria}{\par\smallskip\noindent{\footnotesize\textit{Fonte: elaboração com base no lote experimental da Atividade 3 (2026).}}\par}
\newcommand{\pendencia}[1]{\par\smallskip\noindent{\small\textbf{A completar antes da entrega:} #1}\par}
% Campos que dependem da equipe; editar aqui antes da submissão.
\newcommand{\urlvideo}{\url{https://drive.google.com/drive/folders/1yPalBhHtxjZmD2yZngRQVawkAElJC-Hj?usp=sharing}}
\newcommand{\commitfinal}{\textit{A confirmar após a revisão dos materiais da entrega.}}
\newcommand{\titulotrabalho}{VIRTUALIZAÇÃO, NAMESPACES E CGROUPS: LABORATÓRIO DE CONTÊINERES LINUX}
\newcommand{\integrantes}{Bruno Amancio Ferreira\\Christian Will Silva Santos Nunes\\Iasmin Victoria Teixeira Barreto\\Marcos Vinícius Dantas Aguiar\\Pedro César Figueiredo Carneiro\\Sibele Oliveira Cruz\\Silas Santos da Silva}
\newcommand{\lotecoleta}{\texttt{20261108T013949Z-19921f7f}}
\pgfplotsset{grafico/.style={width=0.47\textwidth,height=4.7cm,
 ymin=0,xmin=0.5,xmax=2.5,axis lines=left,xtick={1,2},
 xticklabels={C1,C2},tick label style={font=\footnotesize},
 label style={font=\footnotesize},scaled y ticks=false,
 grid=major,grid style={gray!18},/pgf/number format/use comma}}

\begin{document}
\hypersetup{pageanchor=false}
\begin{titlepage}
\thispagestyle{empty}
\begin{center}
\textbf{UNIVERSIDADE FEDERAL DE SERGIPE}
\vspace{1.2cm}

\textbf{ATIVIDADE 3 - SISTEMAS OPERACIONAIS (2026.2)}
\vfill
{\bfseries\integrantes}
\vfill
{\large\textbf{\titulotrabalho}}
\vfill
\textbf{SERGIPE}\\\textbf{2026}
\end{center}
\end{titlepage}

\begin{titlepage}
\thispagestyle{empty}
\begin{center}
\integrantes
\vfill
{\large\textbf{\titulotrabalho}}
\vfill
\begin{flushright}
\begin{minipage}{0.55\textwidth}
Relatório técnico apresentado à Universidade Federal de Sergipe como parte da Atividade 3 da disciplina Sistemas Operacionais, período 2026.2.
\end{minipage}
\end{flushright}
\vfill
SERGIPE\\2026
\end{center}
\end{titlepage}

\hypersetup{pageanchor=true}
\pagenumbering{arabic}
\thispagestyle{empty}
\begin{center}\textbf{RESUMO}\end{center}
Este relatório apresenta um laboratório reprodutível de contêineres Linux para observar namespaces e limites de cgroup v2. Scripts próprios, sem o kit de apoio do docente, executaram dois contêineres e uma carga aritmética determinística de uma thread. Foram comparadas quotas de 0,5 e 1 CPU, com três repetições por condição, trabalho fixo de 20.000.000 de iterações, memória de 256 MiB, swap adicional zero e limite de 64 tarefas. Todas as execuções concluíram o trabalho com checksum 2589261568 e limites confirmados nos arquivos do kernel. A mediana do tempo decorrido foi 7,0739 s em 0,5 CPU e 3,3840 s em 1 CPU, razão observada de 2,090 vezes. Os contadores registraram throttling nas três execuções de 0,5 CPU. Os contêineres apresentaram handles distintos de PID, mount e network, embora ambos tivessem PID 1 interno. Em montagem somente leitura, a leitura funcionou e uma tentativa controlada de escrita retornou EROFS. Os resultados distinguem isolamento de visão, controle de recursos e restrições de acesso no kernel compartilhado. As limitações incluem três repetições, interferência não isolada do host e ausência de confirmação dos limites ancestrais externos. Não houve inferência de LLM, uso de GPU ou esgotamento deliberado de recursos.

\vspace{\baselineskip}
\noindent\textbf{Palavras-chave:} Sistemas Operacionais. Contêineres. Namespaces. Cgroups. Virtualização.

\clearpage
\thispagestyle{empty}
\begin{singlespace}\small\tableofcontents\end{singlespace}

\secao{Contribuições, introdução e problema}
\begin{table}[H]
\centering\caption{Participação na apresentação e revisão temática}\label{tab:contribuicoes}
\begin{singlespace}\small
\begin{tabularx}{\textwidth}{>{\raggedright\arraybackslash}p{0.49\textwidth}X}
\toprule\textbf{Integrante} & \textbf{Slides e temas da exposição}\\\midrule
Bruno Amancio Ferreira & 1-2: Abertura, objetivo, organização e percurso do laboratório.\\
Iasmin Victoria Teixeira Barreto & 3-4: Ambiente Windows/WSL2/Docker, diagrama e protocolo experimental.\\
Pedro César Figueiredo Carneiro & 5-6: Namespaces PID, mount e network; numeração dos processos.\\
Christian Will Silva Santos Nunes & 7-8: Carga determinística, checksum, auditoria e demonstração.\\
Marcos Vinícius Dantas Aguiar & 9-10: Limites efetivos do cgroup v2, quotas e throttling.\\
Sibele Oliveira Cruz & 11-12: Seis medições, medianas, gráficos e escopos de memória.\\
Silas Santos da Silva & 13-14: Montagem read-only, menor privilégio, limitações e conclusão.\\
\bottomrule
\end{tabularx}\end{singlespace}
\end{table}
Os nomes são os da equipe da Atividade 2. A Tabela~\ref{tab:contribuicoes} organiza a participação na exposição e a revisão temática. Cada integrante apresenta dois slides consecutivos, com um único apresentador por slide. A integração e a conferência final dos materiais cabem à equipe. O detalhamento está em \path{atividade3/docs/contribuicoes.md}; a assistência de IA está declarada na Seção~\ref{sec:ia}.

A pergunta norteadora é: \textit{como namespaces alteram a visão de processos, mounts e rede, e como cgroups limitam recursos sem transformar um contêiner em uma máquina virtual independente?} O objetivo é responder com comandos, dados e interpretação, conforme o enunciado \cite{enunciado}.

A Atividade 2 forneceu experiência com Linux, coleta, dados sintéticos e análise. Seu pipeline utilizava \texttt{sleep} para simular geração. Nesta atividade, uma carga aritmética de trabalho fixo permite observar quotas de CPU. Os scripts são próprios e não utilizam o kit de apoio do docente.

\noindent\textbf{Repositório:} \url{https://github.com/4rwill/SO_UFS_2026_2_Bruno_Amancio_Christian_Will}.\\
\textbf{Lote principal:} \lotecoleta.\\
\textbf{Vídeo:} \urlvideo\\
\textbf{Commit final da entrega:} \commitfinal

\secao{Ambiente experimental e virtualização}
A coleta principal ocorreu em 07/11/2026, aproximadamente entre 22h39 e 22h40, no fuso America/Sao\_Paulo. Os registros usam UTC: 08/11/2026, entre 01h39 e 01h40. O inventário distingue o host físico, o ambiente Linux efetivo do daemon e o espaço de usuário da imagem.

\begin{table}[H]
\centering\caption{Ambiente efetivamente observado}\label{tab:ambiente}
\begin{singlespace}\small
\begin{tabularx}{\textwidth}{>{\raggedright\arraybackslash}p{0.23\textwidth}>{\raggedright\arraybackslash}X}
\toprule\textbf{Camada} & \textbf{Observação}\\\midrule
Hardware & Intel Core i5-11260H; 6 núcleos, 12 CPUs lógicas; RAM física reportada: 8.358.518.784 bytes.\\
Host & Windows 10 Home Single Language, versão 10.0.19045, build 19045.\\
Linux do daemon & Docker Desktop sobre infraestrutura WSL2; kernel \texttt{6.18.54.1-microsoft-standard-WSL2}.\\
Recursos do daemon & 12 CPUs lógicas; 3.989.606.400 bytes de RAM, aproximadamente 3,72 GiB.\\
Runtime & Docker Desktop 4.96.3; Docker Engine 29.9.2; Compose v5.6.1.\\
Cgroups & v2, driver cgroupfs; controladores cpu, memory e pids observáveis nos contêineres.\\
Imagem & Python 3.12 slim; Debian GNU/Linux 13 (trixie); arquitetura amd64.\\\bottomrule
\end{tabularx}\end{singlespace}
\end{table}
\fontepropria

\begin{figure}[H]\centering
\begin{tikzpicture}[x=1cm,y=0.8cm,every node/.style={font=\small,align=center}]
\draw[gray!65,fill=gray!5] (0,0) rectangle (15.5,4.4);
\node at (7.75,4.02) {Hardware físico e host Windows 10};
\node at (7.75,3.5) {Virtualização WSL2: ambiente Linux intermediário};
\draw[gray!65,fill=white] (0.45,0.35) rectangle (15.05,3.05);
\node at (7.75,2.65) {Kernel Linux compartilhado};
\node at (7.75,2.15) {Daemon Docker Desktop e cgroup v2};
\draw[gray!65,fill=blue!5] (0.85,0.55) rectangle (7.3,1.85);
\draw[gray!65,fill=blue!5] (8.2,0.55) rectangle (14.65,1.85);
\node[font=\footnotesize] at (4.075,1.2) {Contêiner A\\Python e espaço de usuário Debian};
\node[font=\footnotesize] at (11.425,1.2) {Contêiner B\\Python e espaço de usuário Debian};
\end{tikzpicture}
\caption{Camadas do ambiente e kernel compartilhado pelos contêineres.}\label{fig:camadas}
\end{figure}
\fontepropria

Os contêineres compartilham o kernel Linux intermediário; as imagens fornecem bibliotecas e espaço de usuário. O WSL2 estabelece uma camada Linux sobre o Windows. Diferentemente de uma VM com kernel convidado, o contêiner reúne processos sob o kernel compartilhado \cite{namespaces,docker}. Nenhuma VM adicional foi instalada.

A distribuição Ubuntu foi usada na análise e nos testes. O daemon identifica seu ambiente como Docker Desktop; a imagem identifica Debian. \texttt{uname} observa o kernel; \path{/etc/os-release} na imagem descreve seu espaço de usuário.

A imagem oficial foi fixada para todas as observações pelo digest:
\begin{singlespace}\footnotesize\noindent
\path{python@sha256:4840d514c583a4d3d8cad3940e68dee6f0f6f3635c12fcd0a3d74c21a5b93a54}
\end{singlespace}
\noindent{\footnotesize Fontes: \texttt{docker-environment.json} e \texttt{linux-environment.json} no lote principal; inventário complementar em \path{atividade3/results/evidence/host-20261108T014504Z/environment.json}.}

\secao{Namespaces, montagem e menor privilégio}
Dois contêineres independentes da mesma imagem permaneceram ativos durante a coleta. O processo inicial de ambos era Python. Os registros contêm o status de \path{/proc/1}, os handles e campos selecionados de \texttt{inspect}.

\begin{table}[H]\centering
\caption{Comparação dos processos e dos três namespaces}\label{tab:namespaces}
\begin{singlespace}\small
\begin{tabularx}{\textwidth}{>{\raggedright\arraybackslash}Xll}
\toprule\textbf{Recurso} & \textbf{Contêiner A} & \textbf{Contêiner B}\\\midrule
PID no host do daemon & 567 & 612\\
PID interno & 1 & 1\\
PID namespace & \texttt{pid:[4026532292]} & \texttt{pid:[4026532483]}\\
Mount namespace & \texttt{mnt:[4026532287]} & \texttt{mnt:[4026532480]}\\
Network namespace & \texttt{net:[4026532405]} & \texttt{net:[4026532485]}\\\bottomrule
\end{tabularx}\end{singlespace}
\end{table}
\fontepropria
Os handles diferentes sustentam namespaces distintos \cite{namespaces}. O PID 1 interno é relativo a cada namespace; os PIDs 567 e 612, informados pelo Docker, identificam processos distintos no host do daemon. Isso não demonstra isolamento absoluto nem independência de kernel. A rede foi desativada; não houve teste de tráfego.

\noindent{\small Nomes: \path{so-atv3-8b23fdc285-ns-a} e \path{so-atv3-8b23fdc285-ns-b}. Handles e PIDs podem mudar em uma reprodução.}

\subsection{Leitura permitida e escrita negada}
O diretório sintético foi montado em \path{/dados} como somente leitura. A aplicação leu \texttt{texto.txt} e tentou criar exclusivamente \path{/dados/tentativa-escrita.txt}. A Tabela~\ref{tab:mount} apresenta a evidência real.

\begin{table}[H]\centering
\caption{Montagem somente leitura e controles de privilégio}\label{tab:mount}
\begin{singlespace}\small
\begin{tabularx}{\textwidth}{>{\raggedright\arraybackslash}p{0.37\textwidth}X}
\toprule\textbf{Verificação} & \textbf{Observação}\\\midrule
Leitura & Documento sintético recuperado integralmente.\\
Escrita & Negada; errno 30, Read-only file system (EROFS).\\
Mount observado & Opção \texttt{ro} no ponto de montagem \path{/dados}.\\
Identidade & UID 65534, sem root.\\
Capabilities efetivas & \texttt{CapEff = 0000000000000000}.\\
Novos privilégios & \texttt{NoNewPrivs = 1}.\\
Seccomp & Modo 2 observado.\\\bottomrule
\end{tabularx}\end{singlespace}
\end{table}
\fontepropria
Mount namespace define a visão das montagens; a opção \texttt{ro} restringe escrita naquele mount; UID e permissões determinam identidade e acesso; capabilities e no-new-privileges reduzem privilégios \cite{namespaces,dockerrun}. O erro EROFS e a opção \texttt{ro} sustentam a restrição específica da montagem. Não foram montados home, raiz do host, Docker socket ou diretórios sigilosos.

\secao{Configurações, metodologia e métricas}
Um piloto separado de 1.000.000 de iterações levou 0,3095 s sob 0,5 CPU. A calibração escolheu e congelou $N=20.000.000$ antes das seis execuções, respeitando o teto adotado pelo roteiro. O piloto está preservado em \texttt{pilot-0.json} e não entra nas estatísticas principais.

\subsection{Trabalho fixo e correção}
Cada iteração acumula $i^2+3i+7$, módulo $2^{32}$, com $0\leq i<N$. O checksum esperado é calculado independentemente por fórmula fechada:
\begin{equation}
C(N)=\left[\frac{N(N-1)(2N-1)}{6}+\frac{3N(N-1)}{2}+7N\right]\bmod 2^{32}.
\label{eq:checksum}
\end{equation}
Para $N=20.000.000$, o resultado é 2589261568. As seis execuções concluíram todo o trabalho e concordaram com esse oráculo. A carga utiliza uma thread e memória constante, sem LLM, GPU ou espera simulada.

\begin{table}[H]\centering
\caption{Limites efetivos observados antes e depois das cargas}\label{tab:limites}
\begin{singlespace}\small
\begin{tabularx}{\textwidth}{Xllll}
\toprule\textbf{Condição} & \textbf{cpu.max} & \textbf{memory.max} & \textbf{swap.max} & \textbf{pids.max}\\\midrule
C1, 0,5 CPU & 50000 100000 & 268435456 & 0 & 64\\
C2, 1 CPU & 100000 100000 & 268435456 & 0 & 64\\\bottomrule
\end{tabularx}\end{singlespace}
\par\smallskip\noindent{\footnotesize\textit{Fonte: snapshots dos cgroups. ``swap.max'' abrevia memory.swap.max; limites de memória em bytes.}}
\end{table}
Os valores de \texttt{cpu.max} representam tempo agregado por período, em microssegundos. C1 pode consumir 50 ms por período de 100 ms. A quota não fixa afinidade nem reserva um núcleo \cite{cgroups,docker}. A combinação Docker \texttt{memory=256m} e \texttt{memory-swap=256m} corresponde a 256 MiB totais e swap adicional zero, confirmado por \texttt{memory.swap.max}. O limite de 64 tarefas inclui threads e foi confirmado sem provocar esgotamento.

\subsection{Protocolo e escopos de medição}
Foram usados contêineres novos na ordem C1/C2, C1/C2, C1/C2. Imagem, trabalho e demais limites permaneceram iguais. Todos usaram usuário 65534, capabilities removidas, no-new-privileges, rede desativada e raiz somente leitura. Há watchdog interno de 60 s e externo de 90 s para a carga.

O tempo decorrido mede apenas o laço, excluindo criação do contêiner, inicialização do Python e coleta posterior. Tempo de CPU usa \texttt{process\_time}. Trabalho por segundo corresponde a iterações/s, sem equivaler a tokens/s. Os deltas de \texttt{cpu.stat} abrangem snapshots imediatamente antes e depois do laço e pequena sobrecarga de leitura. A análise compara mediana e variação mínimo-máximo de três observações por condição, sem inferência estatística.

\secao{Resultados de CPU e trabalho por segundo}
As seis execuções retornaram status \texttt{completed}, exit code zero e correção verdadeira. O checksum comum foi 2589261568. Os registros completos, o piloto e todos os comandos permanecem no lote \lotecoleta.

\begin{table}[H]\centering
\caption{Seis execuções principais em ordem alternada}\label{tab:seis}
\begin{singlespace}\small
\begin{tabular}{lrrr}
\toprule\textbf{Execução} & \textbf{Wall (s)} & \textbf{CPU do processo (s)} & \textbf{Iterações/s}\\\midrule
C1-1 & 7,0739 & 3,5227 & 2.827.285\\
C2-1 & 3,4952 & 3,4947 & 5.722.124\\
C1-2 & 8,2050 & 4,1136 & 2.437.526\\
C2-2 & 3,3840 & 3,3829 & 5.910.140\\
C1-3 & 6,3852 & 3,1909 & 3.132.258\\
C2-3 & 3,2737 & 3,2732 & 6.109.225\\\bottomrule
\end{tabular}\end{singlespace}
\end{table}
\fontepropria

\begin{figure}[H]\centering
\begin{tikzpicture}\begin{axis}[grafico,ylabel={Tempo decorrido (s)},ymax=10]
\addplot[ybar,bar width=25pt,fill=blue!45,draw=blue!60!black] coordinates {(1,7.073924493) (2,3.384014666)};
\addplot[black,thick,no marks] coordinates {(1,6.385169526) (1,8.205041561)};
\addplot[black,thick,no marks] coordinates {(2,3.273737585) (2,3.495205679)};
\addplot[only marks,mark=*,mark size=1.8pt,black] coordinates {(0.94,7.073924493) (1,8.205041561) (1.06,6.385169526) (1.94,3.495205679) (2,3.384014666) (2.06,3.273737585)};
\end{axis}\end{tikzpicture}\hfill
\begin{tikzpicture}\begin{axis}[grafico,ylabel={Milhões de iterações/s},ymax=7]
\addplot[ybar,bar width=25pt,fill=blue!45,draw=blue!60!black] coordinates {(1,2.827284914) (2,5.910139871)};
\addplot[black,thick,no marks] coordinates {(1,2.437525740) (1,3.132258262)};
\addplot[black,thick,no marks] coordinates {(2,5.722123914) (2,6.109225153)};
\addplot[only marks,mark=*,mark size=1.8pt,black] coordinates {(0.94,2.827284914) (1,2.437525740) (1.06,3.132258262) (1.94,5.722123914) (2,5.910139871) (2.06,6.109225153)};
\end{axis}\end{tikzpicture}
\caption{Mesmo trabalho em C1 (0,5 CPU) e C2 (1 CPU). Barras: mediana; traços verticais: mínimo-máximo; pontos: três observações por condição. Piloto excluído.}\label{fig:cpu}
\end{figure}
\fontepropria

\begin{table}[H]\centering
\caption{Medianas e variação de três repetições por condição}\label{tab:medianas}
\begin{singlespace}\small
\begin{tabularx}{\textwidth}{Xrrrr}
\toprule\textbf{Condição} & \shortstack{\textbf{Wall mediano}\\\textbf{(s)}} & \shortstack{\textbf{Mín.--máx.}\\\textbf{(s)}} & \shortstack{\textbf{CPU mediana}\\\textbf{(s)}} & \shortstack{\textbf{Milhões de}\\\textbf{iterações/s}}\\\midrule
C1 ($n=3$) & 7,0739 & 6,3852--8,2050 & 3,5227 & 2,8273\\
C2 ($n=3$) & 3,3840 & 3,2737--3,4952 & 3,3829 & 5,9101\\\bottomrule
\end{tabularx}\end{singlespace}
\end{table}
\fontepropria
A razão entre medianas de wall, C1/C2, foi 2,090. A redução da mediana foi aproximadamente 52,2\%. É uma observação deste ambiente, com três repetições, sem demonstrar ganho universal ou exatamente igual a dois. A dispersão de C1 foi maior. Ruído do host, escalonamento e variações de frequência ou temperatura são causas possíveis, sem medição que as isole. A alteração do tempo de CPU entre repetições também recomenda cautela.

\secao{Throttling, memória e interpretação}
Os limites efetivos, os contadores e os tempos fornecem evidências complementares. Os deltas positivos de \texttt{nr\_throttled} e \texttt{throttled\_usec} nas três execuções C1 sustentam a ocorrência de restrição de CPU durante o trabalho \cite{cgroups}.

\begin{table}[H]\centering
\caption{Deltas de cpu.stat na janela observada}\label{tab:throttling}
\begin{singlespace}\small
\begin{tabular}{lrrr}
\toprule\textbf{Execução} & \textbf{Usage (s)} & \textbf{nr\_throttled} & \textbf{Throttled (s)}\\\midrule
C1-1 & 3,5247 & 71 & 3,5448\\
C2-1 & 3,4978 & 0 & 0\\
C1-2 & 4,1167 & 82 & 4,0880\\
C2-2 & 3,3845 & 0 & 0\\
C1-3 & 3,1923 & 64 & 3,1918\\
C2-3 & 3,2748 & 0 & 0\\\bottomrule
\end{tabular}\end{singlespace}
\end{table}
\fontepropria
Em C1, a razão CPU/wall ficou próxima de 0,5; em C2, próxima de 1. Nas três execuções C2, os deltas de throttling foram zero nessa janela. A convergência entre limites, contadores e tempos sustenta a interpretação, sem atribuir toda variação à quota ou concluir ausência de toda restrição externa.

\begin{table}[H]\centering
\caption{Escopos distintos das observações de memória}\label{tab:memoria}
\begin{singlespace}\small
\begin{tabular}{lrrr}
\toprule\textbf{Execução} & \shortstack{\textbf{RSS máximo}\\\textbf{do processo (MiB)}} & \shortstack{\textbf{memory.current}\\\textbf{final (MiB)}} & \shortstack{\textbf{memory.peak}\\\textbf{do grupo (MiB)}}\\\midrule
C1-1 & 16,70 & 10,89 & 11,71\\
C2-1 & 16,74 & 11,97 & 11,97\\
C1-2 & 16,77 & 10,89 & 10,89\\
C2-2 & 16,66 & 11,93 & 11,93\\
C1-3 & 16,61 & 10,89 & 11,95\\
C2-3 & 16,78 & 11,92 & 11,92\\\bottomrule
\end{tabular}\end{singlespace}
\end{table}
\fontepropria
RSS máximo vem de \texttt{getrusage} do processo. \texttt{memory.current} representa o uso contabilizado no grupo em dois instantes; a Tabela~\ref{tab:memoria} mostra o final. \texttt{memory.peak} cobre a vida do cgroup, inclusive a inicialização. Compartilhamento e contabilização diferem; não se exige igualdade entre RSS e memória do grupo. Nenhum desses campos é denominado RAM média \cite{cgroups}. Os bytes exatos e as amostras iniciais estão em \texttt{runs.csv}.

O teto de 256 MiB e \texttt{pids.max=64} foram confirmados sem OOM deliberado ou criação excessiva de tarefas. O laboratório demonstra configuração e observação dos limites, sem provar todos os comportamentos ao atingi-los. RAM do cgroup não representa quota de VRAM.

\secao{Reprodutibilidade, validação e uso de IA}
O código está em \path{atividade3/scripts}. O orquestrador \texttt{lab.py} cria snapshots do código e do documento sintético por lote e monta essas cópias read-only. Os contêineres não acessam rede. Identificadores novos preservam resultados anteriores. Comandos, saídas, erros e limites observados ficam registrados.

Commit base da coleta: \path{9f8ad355769af28f7252a5700905476b6dc6a413}, com a Atividade 3 ainda não commitada. O manifesto inclui os hashes dos scripts e os snapshots preservam o código medido. A organização das duas atividades foi publicada em \href{https://github.com/4rwill/SO_UFS_2026_2_Bruno_Amancio_Christian_Will/commit/71626060250925ae7ebbf894991e501057d126d7}{\texttt{7162606}}; o commit final após revisão permanece a confirmar.

A auditoria verificou 68 arquivos, imagem fixa, seis execuções, checksums, ordem pelos horários do daemon, limites antes e depois, contadores, privilégios e limpeza. SHA-256 evidencia integridade interna, sem equivaler a assinatura independente da origem das medições. Da raiz do repositório, executar:
\begin{singlespace}\small\noindent
\texttt{python atividade3/scripts/audit.py}\\
\path{atividade3/results/raw/20261108T013949Z-19921f7f}
\end{singlespace}
As duas linhas acima representam um único comando, com o caminho do lote como argumento.

Seis testes passaram no Ubuntu WSL: oráculo aritmético, limites ausentes ou incorretos, leitura do lote e detecção de adulterações, inclusive checksum ou quota alterados com hashes recalculados. As alterações ocorrem somente em fixtures temporárias.

Tentativas auxiliares no Windows falharam por permissões em temporários; seus logs foram preservados. As seis cargas Docker foram válidas. A análise original utilizou Matplotlib 3.10.8; este relatório incorpora os mesmos valores em gráficos vetoriais PGFPlots.

\subsection{Uso crítico e declarado de IA}\label{sec:ia}
Ferramenta: Codex. O identificador exato do modelo não estava disponível por fonte verificável nesta sessão e deve ser confirmado pela equipe quando possível. A assistência incluiu implementação, coleta autorizada, auditoria, análise, documentação e adaptação do relatório para LaTeX. Não houve uso ou adaptação do kit do docente.

Prompts relevantes: (1) analisar a viabilidade da Atividade 3 com base na Atividade 2; (2) iniciar a implementação sem o kit de apoio; (3) produzir o relatório em LaTeX no formato do relatório da Atividade 2.

Sugestão adotada: carga determinística com fórmula independente, verificada pelo checksum comum e pelos testes. Sugestão corrigida: não tratar o \texttt{/proc} acessível via WSL como equivalente ao host do daemon, após a leitura do PID falhar. Sugestão rejeitada na definição do recorte: utilizar diretamente os tempos do pipeline com \texttt{sleep} como benchmark de CPU. As decisões são sustentadas pelo código, logs e metodologia. A IA não substitui revisão, compreensão e responsabilidade acadêmica da equipe.

\secao{Limitações, conclusão e referências}
A principal ameaça à validade é a ausência de exclusividade e controle do host: há apenas três repetições, alternância fixa e nenhuma medição contínua de carga externa, frequência ou temperatura. A calibração pequena pode sofrer efeitos do período da quota e da inicialização, embora o trabalho principal permaneça fixo e termine com margem para os watchdogs.

Os cgroup namespaces ocultam ancestrais. A tentativa complementar de localizar o PID do daemon no \texttt{/proc} acessível via WSL falhou e foi preservada. Limites ancestrais externos não foram confirmados; RAM e CPUs anunciadas pelo daemon descrevem recursos do ambiente intermediário. A limitação impede afirmar ausência de interferência superior, sem invalidar as leituras locais dos limites.

Em um servidor de inferência compartilhado, quotas de CPU limitam a capacidade de uma carga monopolizar tempo de processamento. Tetos de memória e de tarefas reduzem o impacto de alocações e criação excessiva de processos ou threads. Corpus read-only, usuário não root, capabilities removidas e no-new-privileges restringem ações disponíveis. Permanecem riscos associados ao kernel, runtime, host e permissões concedidas; atualizações e configuração adequada continuam necessárias \cite{namespaces,cgroups,docker}.

Conclui-se que o laboratório evidenciou visões distintas de PID, mount e network, limites efetivos de cgroup v2 e negação de escrita no mount read-only. Com trabalho fixo e checksum comum, a mediana de wall caiu de 7,0739 para 3,3840 s ao elevar a quota de 0,5 para 1 CPU. Os contadores sustentam a ocorrência de throttling em C1. Os resultados se restringem a esta carga aritmética, sem avaliar LLM, GPU, rede, armazenamento sob carga, esgotamento de RAM ou isolamento absoluto entre usuários.

A entrega está prevista para 10/11/2026 e a apresentação para 11/11/2026, com 10--12 minutos de exposição e 3--5 minutos de perguntas. O vídeo está indicado na introdução e em \path{atividade3/VIDEO.md}. A publicação dos materiais no repositório não comprova a submissão no Classroom.

\begingroup
\renewcommand{\section}[2]{}
\begin{singlespace}\footnotesize
\begin{thebibliography}{9}
\bibitem{namespaces} LINUX MAN-PAGES. \textbf{namespaces(7)}. \url{https://man7.org/linux/man-pages/man7/namespaces.7.html}. Consulta: 7 nov. 2026.
\bibitem{cgroups} THE LINUX KERNEL DOCUMENTATION. \textbf{Control Group v2}. \url{https://docs.kernel.org/admin-guide/cgroup-v2.html}. Consulta: 7 nov. 2026.
\bibitem{docker} DOCKER DOCUMENTATION. \textbf{Resource constraints}. \url{https://docs.docker.com/engine/containers/resource_constraints/}. Consulta: 7 nov. 2026.
\bibitem{dockerrun} DOCKER DOCUMENTATION. \textbf{Running containers}. \url{https://docs.docker.com/engine/containers/run/}. Consulta: 7 nov. 2026.
\bibitem{enunciado} UNIVERSIDADE FEDERAL DE SERGIPE. \textbf{Sistemas Operacionais: Atividade 3 - Virtualização, contêineres, namespaces e cgroups}. Enunciado revisado, versão de 5 nov. 2026, disponibilizado pela equipe.
\bibitem{projeto} EQUIPE GENRESEARCH. \textbf{Código, configurações e resultados das Atividades 2 e 3}. 2026. \url{https://github.com/4rwill/SO_UFS_2026_2_Bruno_Amancio_Christian_Will}. Lote da Atividade 3: \lotecoleta. A Atividade 2 é referência de organização e formatação; os dados deste relatório provêm exclusivamente do novo lote.
\end{thebibliography}\end{singlespace}
\endgroup
\end{document}
